% Transcription — fonds Alexandre Grothendieck, shelfmark 154, batch 8 (pages 141-160).
% Established from the facsimile archives/batches/154/batch-08.pdf, cut from a
% copy of 154.pdf fetched through the project's own relay
% (https://grothendieck-relay.onrender.com/source/154.pdf); PDF page k =
% archivists' page 140+k, no cover sheet. Per the skill transcribe-grothendieck.
% The folder is 175 pages in nine batches, transcribed in parallel; this is
% the eighth. Batches 7 and 9 (transcripts/154/) were read first so that the
% notation joins both ends; batches 5-6 fix the rest; otherwise folder 80
% (Systèmes de pseudo-droites).
%
% Pass: Opus 5.5 (claude-opus-5-5), 2026-09-26 — first pass, unchecked against the pages by a human.
%
% Dating and title. \dating is the folder's own date, copied verbatim from
% src/content/catalogue.ts; the folder belongs to the inventory group
% « Espaces stratifiés » (150-151), but since it has a date of its own the
% group's range is not added. Title from the catalogue, trailing full stop
% dropped; « [Système de pseudo-droites] » is bracketed there because the
% archivists supplied it.
%
% Paper. Every leaf is fan-fold computer listing paper. Printed banners and
% logs (a plotter log dated 25/8/82 on 141; « DATASETS PICTURES » tables on
% 144, 156; JES2 logs and notices of August 1982 on 145, 150, 151, 154, 155,
% 158, 160, several printed upside down) are machine dates of the paper, not
% his, and are not transcribed.
%
% Pages skipped: none. Every leaf carries drawings or text of his.
% Pages summarised: none. Pages given only as a figure described in a French
% \note, with their few labels set: 142, 143, 146, 147, 148, 149, 151, 158,
% 160; figures with a formula, a list or a legend: 141 (a count ~ n^4/8 and
% the list « p.r. stables / sous-stables / critiques »), 144 (the sequence of
% permutations 1'-9' of the pseudo-lines A...G), 145 (a braid word in the
% sigma_i), 155 (two legends).
% Pages transcribed from his hand: 150, 152, 153, 154, 156, 157, 159.
%
% His pagination and dates. No page number of his. One date in his hand:
% « 1.1.84 », underlined, at the head of p. 159 (the last digit is written
% as his open 4). Numbered runs: 1'-9' on p. 144; lines 1-12 and 12-19 of the
% braid diagrams and the path 1-19 on p. 145; vertices 1-19 and « 17 = 1 » on
% p. 147; the types of spindles 1) 3) 3) on p. 152 (the second number
% reads « 3) » as written) and the axioms a)-d) of a « fuseau » there;
% 1) on p. 153; (I_1), (I_2), (I_3), (I_n) (n = 7) on p. 154; a)-d) on
% p. 159. p. 141 does not visibly continue p. 140; p. 160's drawings (the
% cylinder facets F, F', bands a a', b b') illustrate p. 159 and lead into
% the drawings that open batch 9 (p. 161: discs with arcs between a and b,
% rim labels 1-5, 1'-5').
%
% What the batch is about. Continuing batch 7 (spindles, « repères » and the
% involutions sigma, sigma_0, sigma_1, sigma_2 on them):
%   141-149  counting (1/2 N(N-1) - n(n-1)(n-2)/2 ~ n^4/8, N = n(n-1)/2);
%            sweeps of a disc of seven pseudo-lines A...G, recorded as
%            sequences of permutations (p. 144) and braid diagrams with
%            generator words sigma_i (p. 145).
%   150      The cellular surface of relative positions X* = Rel(Sigma) of
%            (X, Sigma): at a vertex rho (critical p.r. D_rho, S_rho =
%            D_rho ∩ S(Sigma)) the edges leaving rho correspond to pairs
%            (s, omega), s in S_rho, omega an orientation; the edge a_s,omega
%            and the « opposite » edge b_s,omega between stable p.r.; the
%            difficulty is the local structure near an edge a_L (semi-stable
%            p.r. L) -- which critical position is its endpoint.
%   152      Types of spindles: couloirs (both banks supercritical),
%            semi-couloirs, ordinary spindles; a « fuseau » is a connected
%            component C cut out by two pseudo-lines D, D' with a)-d); each
%            spindle determines a sub-stable p.r. and conversely.
%   153      For a non-principal p.r. D with card S_D >= 2: is every principal
%            p.r. through S_D a specialisation of D?  The stable generisations
%            of D are the vertices of a nu-cube (2^nu vertices, nu 2^(nu-1)
%            edges), the sub-stable ones its edges.
%   154-155  Elementary sweeps (I_1), (I_2), (I_3), (I_n) of relative
%            positions, symmetric or asymmetric.
%   156-157  The operations on (D_0, beta*, omega*): sigma, sigma_0, sigma_2
%            change the signs of beta*, omega*; rho_s = sigma_2 sigma_1,
%            rho_f = sigma_1 sigma_0, rho_s rho_f = sigma; the dual real
%            projective planes X, X-check and orientations (alpha, beta,
%            omega « associés » vs alpha', beta', omega' « non associés »).
%   159      (1.1.84, whole page cancelled by a green cross) Exceptional
%            facets of (X, K): phi^{-1}(X_1) = L ∪ K_0; the structure near
%            exceptional edges; the privileged bank; the p.r. attached to an
%            exceptional face.
%
% Notation fixed once for the batch, following batch 7: \varpi for his
% barred omega, \omega for an orientation of D, \sigma, \sigma_0, \sigma_1,
% \sigma_2 for his barred sigmas, « p.r. » for position relative; X^{*} for
% the surface of positions on p. 150 (his exponent is a star), \mathcal{X}
% for his cursive X on p. 159 (as in batch 9); \check{X} for the dual plane.
%
% Readings to check first: p. 150 (the struck lines after (D_rho,s,omega),
% « très écartée », the last paragraph); the NB and the first alinea of
% p. 152; p. 153, second paragraph; p. 157, the whole text; p. 159 b), the
% paragraph under the drawing and d); the braid word on p. 145 from line 9.

\documentclass[11pt,a4paper]{article}
% Le préambule commun à toutes les transcriptions du fonds Grothendieck.
%
% Il tient en une idée : **l'incertitude se compose, elle ne se résout pas.**
% Une transcription qui lisse un mot douteux a détruit la seule chose pour
% laquelle on la faisait — un lecteur doit pouvoir savoir, à chaque endroit,
% ce qui était sur le feuillet et ce qui a été ajouté. Les six macros qui
% suivent sont donc l'appareil critique, et non de la décoration.
%
% Les mêmes commandes sont reconnues par `scripts/render.mjs`, qui produit la
% vue de lecture du site. Une macro ajoutée ici doit l'être aussi là-bas,
% faute de quoi le rendu échouera bruyamment — ce qui est voulu : un rendu
% silencieusement faux est pire qu'un rendu absent.

\ProvidesPackage{grothendieck}[2026/08/08 Transcriptions du fonds Grothendieck]

\usepackage{amsmath,amssymb,amsthm}
\usepackage{mathrsfs}
\usepackage{stmaryrd}
% Les diagrammes commutatifs sont partout dans ce fonds : c'est la langue
% même de la géométrie algébrique de Grothendieck, et une transcription qui
% les rendrait en prose ne transcrirait rien.
\usepackage{tikz-cd}
% Français par défaut — c'est la langue des feuillets — mais l'anglais est
% chargé aussi : la traduction et le résumé sont des documents anglais, et la
% typographie française (espaces avant les deux-points) y serait une faute.
% Ces documents basculent par \selectlanguage{english} en tête de corps.
\usepackage[english,french]{babel}
\usepackage{fontspec}
\usepackage{geometry}
\geometry{a4paper,margin=25mm,marginparwidth=28mm}
\usepackage{marginnote}
\usepackage{xcolor}
% ulem, not soul. `soul`'s \st and \ul are built on a character-by-character
% scan that XeLaTeX's font machinery breaks: the document fails with an
% undefined control sequence rather than degrading. ulem does the same two jobs
% and is Unicode-engine safe. `normalem` stops it hijacking \emph, which the
% transcriptions use for Grothendieck's own emphasis.
\usepackage[normalem]{ulem}
\usepackage{hyperref}
\hypersetup{colorlinks=true,linkcolor=black,urlcolor=black,pdfborder={0 0 0}}

% --- Métadonnées du lot -----------------------------------------------------
% Reprises telles quelles par le script de rendu, qui les affiche en tête de
% la vue de lecture. Elles fixent ce que le fichier prétend couvrir : sans
% elles, une transcription flotte sans point d'attache dans un fonds de seize
% mille feuillets.

\newcommand{\folder}[1]{\gdef\grothFolder{#1}}
\newcommand{\batch}[1]{\gdef\grothBatch{#1}}
\newcommand{\leaves}[2]{\gdef\grothFirst{#1}\gdef\grothLast{#2}}
\newcommand{\foldertitle}[1]{\gdef\grothFoldertitle{#1}}
\gdef\grothFolder{?}\gdef\grothBatch{?}\gdef\grothFirst{?}\gdef\grothLast{?}\gdef\grothFoldertitle{}

% --- Le résumé --------------------------------------------------------------
%
% Il ouvre la lecture modernisée, sous le titre « Résumé », et il est composé
% en petit. Ce n'est pas un ornement : c'est le seul passage du document qui
% ne soit pas des mathématiques, et un lecteur doit pouvoir le reconnaître
% avant de l'avoir lu — puis le sauter s'il connaît déjà le sujet, ou s'y
% arrêter s'il ne le connaît pas. Le filet qui le suit marque la reprise du
% corps.

\newenvironment{resume}
  {\small\setlength{\parskip}{0.45em}}
  {\par\medskip\hrule\medskip}

% --- Mots-clés ---------------------------------------------------------------
%
% En anglais, en clôture du résumé de la lecture modernisée : le vocabulaire
% moderne sous lequel chercher ce que les feuillets construisent. Le manifeste
% du site (`npm run manifest`) les extrait pour étiqueter la cote entière —
% cette ligne est la source unique des tags, il n'y a pas de fichier à côté.
\newcommand{\keywords}[1]{\par\smallskip\noindent
  {\footnotesize\textsc{Keywords} --- \textit{#1}}\par}

% --- L'appareil critique ----------------------------------------------------

% Le numéro de feuillet, en marge. C'est l'ancre qui permet de régler un
% désaccord sur une formule en regardant *un* feuillet plutôt qu'une chemise
% de six cents. Le site s'en sert aussi pour tourner les pages du fac-similé
% quand on fait défiler la transcription.
\newcommand{\leaf}[1]{%
  \marginnote{\footnotesize\color{gray!70}\textsf{#1}}%
  \label{leaf:#1}\ignorespaces}

% Illisible : on ne devine pas. Un mot inventé qui se lit comme les autres est
% le pire résultat possible.
\newcommand{\ill}{\textcolor{red!60!black}{[\,\ldots\,]}}

% Lecture douteuse : le mot est donné, et signalé comme incertain.
\newcommand{\uncertain}[1]{\uline{#1}}

% Ajout de l'éditeur — une lettre manquante, un mot sous-entendu.
\newcommand{\add}[1]{\textcolor{blue!50!black}{[#1]}}

% Biffé par Grothendieck lui-même. Ce qu'il a rayé fait partie du document :
% une rature dit souvent où la pensée a changé de direction.
\newcommand{\struck}[1]{\sout{#1}}

% Note du transcripteur, sur l'état matériel du feuillet.
\newcommand{\note}[1]{\par\smallskip{\small\color{gray!60!black}\textit{[#1]}}\par\smallskip}

% Note marginale de Grothendieck, distincte du corps du texte.
\newcommand{\marginal}[1]{\par\smallskip\noindent\hspace{1em}%
  {\small\color{gray!60!black}#1}\par\smallskip}

% --- Titre ------------------------------------------------------------------

\AtBeginDocument{%
  \begin{center}
    {\large\bfseries Fonds Alexandre Grothendieck --- cote n\textsuperscript{o}~\grothFolder}\\[2pt]
    {\itshape\grothFoldertitle}\\[4pt]
    {\small feuillets \grothFirst--\grothLast\ (lot \grothBatch)}\\[2pt]
    {\footnotesize\color{gray!60!black}Transcription établie d'après le fac-similé numérisé
     par l'Université de Montpellier.\\
     Les lectures incertaines sont soulignées, les illisibles notées [\,\ldots\,],
     les ajouts entre crochets.}
  \end{center}
  \vspace{1em}}

\endinput


\folder{154}
\batch{8}
\pages{141}{160}
\dating{1983-1984}
\watermark{Édition de démonstration}
\foldertitle{[Système de pseudo-droites] : notes manuscrites (1983-1984, s.d.)}

\begin{document}

\page{141}
\note{page de dessins au crayon et à l'encre : deux fuseaux allongés (bigones) traversés de pseudo-droites, l'un en haut à gauche, l'autre en bas à droite ; au milieu, une longue droite marquée $D$ et $D'$ (plusieurs fois), coupée de croisements, un chemin brisé en zigzag qui la longe, des hachures le long de la droite ; à gauche, un disque traversé d'arcs qui se croisent (une sorte de fuseau à sommets marqués) et un grand disque traversé de pseudo-droites ; en haut à droite, des arrangements de droites au crayon ; un rectangle au crayon}

$(n-1)(n-2)$
\note{écrit au crayon sous le disque de gauche, souligné, avec un « $n$ » à gauche au-dessous et un « $2$ » au-dessous}
\[
\frac{1}{2}\,\frac{n(n-1)}{2}\left(\frac{n(n-1)}{2} - 1\right) - \frac{n(n-1)(n-2)}{2} = {}\struck{\ill{}} \sim \frac{1}{8}\,n^4
\]

p.r.\ stable

p.r.\ sous-stables

p.r.\ \emph{critiques} (ou \uncertain{minimales} \uncertain{pour} la \uncertain{spécialisation})
\note{en bas à droite, au crayon, sous le second fuseau}

\page{142}
\note{page de dessins à l'encre. En haut à gauche, un rectangle partagé par deux droites horizontales, les colonnes marquées en haut $A\,B\,C\,D\,E\,F\,G\,H$ (flèche vers la droite) et en bas $H\,G\,F\,E\,D\,C\,B\,A$ (lecture de l'avant-dernière lettre incertaine), une flèche vers le bas à gauche. Au milieu, un arbre binaire de flèches issu de $D\,E\,F\,G$ ; un fuseau allongé traversé de droites ; un croisement de deux droites sur deux horizontales. En haut à droite, un diagramme en « fils » : sept fils verticaux $A, B, C, D, E, F, G$ qui se croisent deux à deux, sur un réseau de lignes horizontales numérotées à droite $1$ à $11$ et à gauche $1'$ à $9'$ ; en bas, l'ordre d'arrivée $G\,A\,E\,D\,C\,F\,B$. En bas à droite, le même principe sur un quadrillage, colonnes $A$ à $G$, lignes $1$ à $11$, l'arrivée marquée $G\,A$. En bas au milieu, un cylindre dont le bord supérieur porte $G, F, E, D$ et $A, B, C$, le bord inférieur $D, E, F$ et $B, A, G$}

\page{143}
\note{croquis au crayon et à l'encre, sans texte : des arrangements de trois ou quatre droites qui se coupent, certaines en étoile autour d'un point ; à gauche et à droite, des chemins en tirets (bruns, verts) qui longent un segment épais entre deux croisements}

\page{144}
\note{à gauche, un arbre de flèches : $D, E, F, G$ en haut, puis les couples $DE$, $FG$, $DG$, $EG$, $CG$, $BG$, $AG$, les flèches descendant jusqu'à une ligne de six flèches ; chaque couple est relié par un trait à la ligne correspondante de la colonne de droite, où le couple qui se croise est souligné d'une accolade}
\[
\begin{array}{ll}
\struck{A\,B\,C\,D\,E\,F\,G} & \\
A\,B\,C\,D\,E\,F\,G \quad \struck{A\,B\,C\,E\,D\,G\,F} & 1' \\
A\,B\,C\,E\,D\,G\,F & 2' \\
A\,B\,C\,E\,G\,D\,F & 3' \\
A\,B\,C\,G\,E\,D\,F & 4' \\
A\,B\,G\,C\,E\,D\,F & 5' \\
A\,G\,B\,C\,E\,D\,F & 6' \\
(A\,G)\,C\,B\,E\,D\,F & 7' \\
(A\,G)\,C\,E\,B\,D\,F & 8' \\
(A\,G)\,E\,C\,D\,B\,F & 9' \\
G\,A\,E\,D\,C\,F\,B &
\end{array}
\]
\note{les numéros $1'$ à $9'$ sont écrits à droite, un peu décalés vers le bas ; les couples lus $(A\,G)$ des lignes $7'$ à $9'$ sont entourés, le premier surchargé ; les couples qui vont se croiser sont soulignés ; dans la dernière ligne, les lettres $D$, $C$, $F$ sont surchargées. À droite, un diagramme de fils sur des lignes numérotées $5'$ à $11'$ ; en bas, trois arrangements de droites en réseau losangé, l'un traversé de lignes en tirets ; en haut à droite, une droite coupée de deux autres}

\page{145}
\note{grand dessin à l'encre et au crayon : un disque dont le bord porte les extrémités de sept pseudo-droites $A, B, C, D, E, F, G$ (chacune deux fois), un faisceau de fines courbes au crayon qui joignent deux points opposés du bord (un balayage), et un chemin de flèches brunes numérotées $1$ à $19$ qui passe de croisement en croisement ; les deux extrémités $E$, $D$, $C$ du côté droit sont fléchées. En haut à droite, des croquis de droites. En bas, deux diagrammes de fils : à gauche, les fils partent dans l'ordre $D\,F\,B\,E\,G\,C\,A$ et arrivent dans l'ordre $G\,F\,E\,D\,C\,B\,A$, sur des lignes numérotées $12$ à $19$ ; à droite, les fils $A\,B\,C\,D\,E\,F\,G$ (numérotés au-dessus $1$ à $6$ entre les fils) sur des lignes numérotées $1$ à $12$, une ligne en tirets entre $4$ et $5$}

$2 \quad 5 \quad 7 \quad 3 \quad 3 \quad 2$
\note{au-dessus du diagramme de droite, entre deux traits}

\[
\sigma_3,\ \sigma_4,\ \sigma_5,\ \sigma_2,\ \sigma_6,\ \sigma_1,\ \sigma_5,\ \sigma_3,\ \sigma_6,\ \ill{},\ \sigma_3\ \sigma_5,\ \sigma_2\ \sigma_6,\ \ill{}
\]
\note{colonne des générateurs, à droite du diagramme de fils de droite, en face des lignes $1$ à $12$ ($\sigma_2, \sigma_6$ en face des lignes $4$ et $5$) : à chaque ligne, le croisement entre les fils de rangs $i$ et $i+1$ est noté $\sigma_i$. À partir de la ligne $9$, la colonne est surchargée et en partie noyée dans une tache d'encre ; le neuvième terme, lu $\sigma_6$, et les derniers, écrits en deux colonnes, sont d'une lecture incertaine. Le diagramme de gauche porte une colonne semblable, illisible}

\page{146}
\note{quatre disques au crayon et à l'encre, sans texte : un disque traversé de pseudo-droites qui se croisent ; un cercle vide ; un disque dont deux points opposés du bord sont joints par un faisceau dense de courbes brunes ; un disque traversé de courbes enchevêtrées ; un disque traversé de quelques droites entre deux points marqués du bord}

\page{147}
\note{grand dessin à l'encre et au crayon : un disque aux deux pôles marqués (en haut et en bas), joints par un faisceau de fines courbes (un balayage), traversé de pseudo-droites ; un réseau de flèches (noires et brunes) va de croisement en croisement, les sommets numérotés de $1$ à $19$ ; plusieurs numéros sont groupés et entourés : « $1, 2, 3, 4$ » en haut à gauche, « $18, 19$ » (deux fois), « $15, 16$ », « $10, 11$ », « $12, 20$ » (lecture incertaine) ; en haut à droite, « $17 = 1$ ». Des régions hachurées au bord. En bas à droite, un petit disque au crayon rempli d'un réseau de flèches ; en haut à droite, un croisement de trois droites sur une droite hachurée}

\page{148}
\note{croquis au crayon et à l'encre, sans texte : un arrangement de cinq droites (en étoile) ; un disque aux deux pôles marqués, joints par un faisceau de courbes, traversé de pseudo-droites aux croisements marqués, des coins hachurés ; un réseau de flèches au crayon ; en bas, une droite en tirets qui passe par des croisements de droites, un croisement marqué d'un gros point et d'une petite flèche de rotation}

\page{149}
\note{croquis au crayon et à l'encre, sans texte : en haut, deux fuseaux allongés enfilés sur une droite ; au milieu, une droite brune passant par trois croisements en étoile, et un fuseau ; en bas, un arrangement de droites dont deux en tirets (vert et brun), trois croisements marqués ; un petit graphe en losange dans un ovale ; des segments isolés}

\page{150}
$\underline{\Sigma} = (X, \Sigma)$, d'où surface cellulaire des positions relatives $X^{*} = \uncertain{\mathrm{Rel}}(\underline{\Sigma})$.
\note{le nom du foncteur, lu « Rel », est surchargé ; l'exposant de $X$ est une étoile, ici et dans toute la page}

Soit un sommet $\rho$ \add{de} $X^{*}$, correspondant à une position relative critique, \struck{A} $D_\rho$, soit $S_\rho = D_\rho \cap S(\Sigma)$. (Une orientation \struck{$\varpi$} de $X^{*}$ en $\rho$ correspond à \struck{\ill{}} une orientation $\varpi_\rho$ de $D_\rho$.)
\note{la phrase entre parenthèses est écrite dans l'interligne, sous la précédente ; « $\rho$ » est ajouté au-dessus de « sommet », « $\varpi_\rho$ » au-dessus de la fin de la ligne}

Alors les \struck{\ill{}} arêtes sortant de $\rho$ \add{de} $X^{*}$ correspondent aux \struck{\ill{}} éléments \add{$(s, \omega)$} de $S_\rho$, munis d'une orientation, [soit $\vec{a}_{s,\omega}$ l'arête]. À un tel couple $(s, \omega)$ correspond la p.r.\ sous-stable $(D_{\rho,s,\omega})$\struck{, \ill{}}, \struck{qui pourra être \ill{} \ill{} \ill{}} \struck{\ill{} la \uncertain{génér.}\ $(\rho, \ill{}\,\omega)$ correspondante : p.r.\ stable} $\to$ à l'arête $a_{s,\omega}$ sous-jacente à $\vec{a}_{s,\omega}$ ; soit $b_{s,\omega}$ l'arête \uncertain{opposée}, orientée de façon que $(\vec{a}_{s,\omega}, \vec{b}_{s,\omega})$ soit \uncertain{directe} pour $\varpi$, donc $\vec{b}_{s,\omega}$ va d'une p.r.\ \emph{stable} vers une autre --- c'est le passage de la générisation de $D_{s,\omega}$ en $s$ « \uncertain{très} \uncertain{écartée} » vers la gén.\ \ill{} \struck{\uncertain{gauche}} \add{droite}.
\note{« $(s, \omega)$ » est écrit au-dessus de « éléments » ; « soit $\vec{a}_{s,\omega}$ l'arête » est ajouté au-dessus de la ligne et renvoyé par une accolade ; les deux lignes biffées qui suivent $(D_{\rho,s,\omega})$ sont barrées de deux grands traits obliques ; « droite » est écrit au-dessus du dernier mot biffé}

\note{à gauche, un dessin : une droite fléchée $D_\rho$ coupée d'autres droites en un sommet $s$, l'orientation $\omega$ marquée par une petite flèche de rotation, et deux chemins en tirets (brun et vert) qui la longent, l'un fléché et marqué $D_{\rho,s,\omega}$ ; au-dessous, $\vec{\omega}_{D_\rho}$. À droite, un petit dessin : au sommet $\rho = \rho_D$, une flèche épaisse $\vec{a}_{s,\omega} = \vec{a}$, l'orientation $\varpi = \varpi_\rho$ et un contour en tirets fléché ; plus à droite, une arête $a = a_L$ entre deux faces $F_{\Delta'}$ et $F_{\Delta}$, avec les orientations $\omega_a$ et $\varpi_a$ ; au-dessous, une droite $L$ fléchée ($\omega_L$) passant par un sommet où est marquée l'orientation $\omega_s$}

La difficulté conceptuelle n'est pas dans l'étude locale de $X^{*}$ au voisinage de $\rho$, i.e.\ au voisinage de $D$, mais dans l'étude locale de $X^{*}$ au voisinage d'une face, correspondant à une p.r.\ stable $\Delta$, mais dans la description de $X^{*}$ au voisinage d'une \emph{arête} \add{$a_L$}, correspondant à une position r.\ semi-stable donnée, soit $L$. Il y a un peu de difficulté : \uncertain{voir} les faces incidentes \struck{la face $F$} à $a_L$, et leur circuit des arêtes (contenant $a_L$), mais pas \ill{} : \ill{} les deux \add{(?)} sommets incidents à $a_L$.
\note{les deux « mais » sont lus tels qu'ils sont écrits ; « $a_L$ » est ajouté au-dessus de « arête », « (?) » au-dessus de « deux » ; la phrase continue en haut de la colonne de droite}

Notons que la donnée d'une orientation superficielle le long de $a_L$ revient à celle d'une orientation de $L$, celle d'une orientation de $a_L$ à celle d'une orientation superficielle en le \struck{pt} sommet $s \in S_L$, donc une \emph{rive} de $L$\struck{\ill{}} en $s$ revient à celle d'une rive de $a$, i.e.\ d'une des faces incidentes.

Il s'agirait donc, en termes de $(L, \omega_s)$, de \ill{} \ill{} le sommet extrémité de $\vec{a} = (a, \omega_a)$, \struck{sans doute} \uncertain{si} possible en termes d'une position relative critique, déduite de la p.r.\ « \uncertain{bloquée} » autour de $s$, dans le sens indiqué par l'orientation $\omega_s$, i.e.\ par la rive correspondante de $L - \lbrace s \rbrace$. L'ennui, c'est qu'il y a plusieurs choix naturels qui viennent à l'esprit, et qu'il n'est pas sûr qu'aucun de ces \ill{} \uncertain{ne} \uncertain{marche} !

\page{151}
\note{page de dessins à l'encre et au crayon, sans texte : une vingtaine de disques traversés de pseudo-droites (brunes, noires, vertes), certains arcs du bord renforcés en vert ; dans plusieurs disques, les régions sont marquées alternativement de deux signes, l'un en forme de $\varepsilon$, l'autre en forme de chiffre $3$ à longue queue (un coloriage en deux classes, semble-t-il) ; dans un disque, une étoile de trois droites jaunes et noires, deux d'entre elles hachurées, un triangle central hachuré ; un disque divisé en secteurs par huit rayons ; des pentagones ; de petits disques verts ; deux longues droites au crayon traversent la page en croix}

\page{152}
\note{en haut, deux dessins de fuseaux : l'un au crayon, traversé de pseudo-droites, avec deux flèches noires d'un bord à l'autre ; l'autre en couleur, bordé d'orange, rempli de courbes vertes et de tirets bruns qui joignent les deux pointes, traversé d'un réseau de droites brunes qui se croisent en losanges, avec de petites flèches aux croisements}

\emph{NB} Même une fois complété par les fuseaux adjacents, le \ill{} en flèches \ill{} \ill{} \ill{} (car il y a \ill{} \ill{} connexion.
\note{colonne de droite, sous le dessin en couleur ; « NB » est souligné en brun ; la lecture de ce passage est très incertaine}

Types de fuseaux :

1) couloirs : les deux rives \struck{sont} \add{bords} des fuseaux supercritiques (ou fuseaux \struck{\uncertain{bi-supercritiques}})

\uncertain{3)} semi-couloirs : une et une seule des deux rives est supercritique ;

3) fuseaux \emph{ordinaires} \struck{\ill{}}

\uncertain{f.\ terminaux} (\ill{} \struck{ou f.\ critiques} \ill{})
\note{« bords » est écrit au-dessus du mot biffé ; le numéro du deuxième alinéa se lit « 3) », comme celui du troisième ; au-dessous du troisième, un mot biffé, lu « ordinaires ». Les mots « \uncertain{f.\ terminaux} » et la parenthèse qui suit sont écrits à gauche, sous « Types de fuseaux », reliés par une grande courbe à la parenthèse du premier alinéa}

\struck{ordinaire}

Fuseau : comp.\ connexe $C$ déterminée par deux ps.\ droites $D, D'$ (qui peuvent être $\in \Sigma$) telles que \struck{1) $D \neq D'$} :

a) $\Sigma \cup \lbrace D, D' \rbrace = \Sigma'$ syst.\ de ps.\ dr.

b) $D \neq D'$ \struck{\ill{}} et \ill{} $= D \cap D'$ sont des sommets de $\Sigma$

c) $C$ est un couloir de $\Sigma'$ \struck{\ill{}} !

d) $D$ et $D'$ sont critiques pour $\Sigma$
\note{« $= \Sigma'$ » est écrit au-dessus de la ligne a), rattaché par une accolade à « syst.\ de ps.\ dr. » ; le « $C$ » de la première ligne est écrit dans la marge droite}

d) signifie aussi que \struck{tt triangle} de contraction sur la « rive \emph{libre} » de $D$ ou $D'$ est réduit à un sommet, et qu'à l'intérieur du fuseau, chaque sommet \struck{\ill{}} d'une rive est relié à un sommet sur l'autre (cas $\Sigma$ de type I ou II), sauf dans les cas triviaux, un fuseau est déterminé par les (classes d'isotopie) de ses rives.
\note{colonne de gauche, sous le troisième dessin ; au-dessus de « triangle » biffé, un mot illisible. Le troisième dessin, au crayon : un fuseau entre les deux pseudo-droites $D$ et $D'$, coupé de droites, et un chemin brun en zigzag qui va d'une rive à l'autre de croisement en croisement}

Chaque fuseau détermine une p.r.\ sous-stable, et inversement. Corr.\ $1$-$1$ entre ((classes d'isotop.\ des) \uncertain{fus.}) dans $X$ et p.r.\ sous-stables).
\note{colonne de droite, en bas ; la parenthèse intérieure est écrite au-dessus de la ligne}

\page{153}
\note{en haut, au crayon, deux dessins : un fuseau dont le sommet de droite est traversé de droites marquées $\Delta_1$, $\Delta_\nu$, $\Delta_3$ ; un chapelet de deux fuseaux sur une droite, deux flèches brunes au sommet commun et au sommet de droite}

Soit $D$ position relative non principale, $S_D$ \add{$= D \cap S$} (\struck{l'ens.}) des sommets multiples de $D$ ($\operatorname{card} S_D \geqslant 2$).

1) Toute position relative principale \add{(donc critique)} qui passe par $S_D$ est-elle une spécialisation de $D$ ?

Pour celles \add{$D' = D'_S$} qui (sont des spécialisations de $D$, \uncertain{i.e.}\ elles ne sont pas supercritiques), à chacune d'elles est associé un syst.\ de flèches d'incidence \uncertain{aux} $\Delta_s$ ($s \in S_{D'} \setminus S_D$) qui « passent » par $s$, et ne sont pas parallèles à \ill{} ($s \notin S_D$). Pour un $S$ donné, un système de flèches \ill{} détermine \ill{} $D$.
\[
\prod_{s \in S_{D'} \setminus S_D} \mathcal{E}_{D',s}
\]
\note{les additions « $= D \cap S$ », « (donc critique) » et « $D' = D'_S$ » sont écrites au-dessus des lignes ; la lettre lue $\mathcal{E}$ est surchargée}

Il y a aussi le système des \struck{\ill{}} \add{p.r.\ stables} qui généralisent $D$,

Les générisations stables de $D$ correspondent aux sommets d'un cube \struck{\ill{}} de dimension $\nu$, dont les arêtes sont les générisations sous-stables, et la relation d'incidence entre sommets et arêtes est la relation de spécialis\supplied{ation}.

Produit des diagonales d'un polygone \add{convexe} \ill{} \struck{de degré} $2\nu$

$2^{\nu}$ sommets, $\nu 2^{\nu - 1}$ arêtes
\note{colonne de droite ; l'exposant du premier $2$ est surchargé. En bas, trois dessins : une étoile de trois droites au crayon ; la projection d'un cube (hexagone à sommets marqués) entourée de deux ellipses brunes ; un cercle portant huit petites marques}

\page{154}
\note{en haut, trois esquisses au crayon très pâles d'une bande qui s'élargit (sans étiquettes)}

balayages élémentaires des positions relatives

\note{trois dessins, marqués au crayon $(\mathrm{I}_1)$, $(\mathrm{I}_2)$, $(\mathrm{I}_3)$ : une droite verte horizontale (la position de départ), bordée de deux courbes brunes numérotées $1$ (au-dessous) et $2$ (au-dessus), entre deux traits verticaux en tirets ; au crayon, des faisceaux de droites de $\Sigma$ passant par des sommets marqués. Dans $(\mathrm{I}_1)$, la courbe $2$ fait une bosse en tirets pour passer par le sommet $a$ ; dans $(\mathrm{I}_2)$, les deux courbes s'écartent en un fuseau, $2$ passant par $a$ en haut et $1$ par $b$ en bas, une flèche double brune de $b$ à $a$ ; dans $(\mathrm{I}_3)$, $2$ passe par $a$ et $b$, $1$ par $c$, un chemin de flèches brunes $a$, $c$, $b$. Chaque dessin est suivi d'une accolade et d'un texte}

$(\mathrm{I}_1)$ pas d'autre sommet de $\Sigma$ ici que $a$, ni d'autres traces des droites de $\Sigma$ que \struck{\ill{}} celles marquées (passant par $a$).

\emph{asymétriques}, la position $1$, n'est pas \struck{\uncertain{générique}} p.r.\ critique

$(\mathrm{I}_2)$ pas d'autres sommets de $\Sigma$ que $a$ et $b$, ni d'autres traces des droites de $\Sigma$ que celles marquées (passant par $a, b$)

\emph{symétriques}, la position $1$ est critique ssi la position $2$ l'est

$(\mathrm{I}_3)$ pas d'autres sommets de $\Sigma$ ici que $a, b, c$, ni d'autres traces de droites de $\Sigma$ que celles marquées (passant par $a, b, c$))

\emph{asymétriques}, néanmoins $1$ est critique ssi $2$ l'est

\note{sous les trois textes, une grande accolade qui les réunit, puis :}
balayages élémentaires des p.r.\ critiques à une \ill{}
\note{le premier mot de « élémentaires » est surchargé}

$(\mathrm{I}_n)$ \quad $(n = 7)$

Cas asymétrique car $n$ impair
\note{à gauche, le dessin correspondant : la même bande, les deux courbes brunes passant alternativement par sept sommets, quatre en haut et trois en bas (lecture du nombre incertaine), reliés par des faisceaux de droites au crayon qui dessinent des zigzags}

\page{155}
\note{page de dessins au crayon et à l'encre brune et verte, dans les conventions de la page 154 : en haut, un fuseau brun bordé de droites qui se croisent en losanges, et au-dessous un fuseau au crayon ; au milieu, des croisements de droites, un fuseau au crayon parcouru d'un zigzag brun entre ses deux rives (sept sommets) ; une courbe brune qui s'abaisse en tirets pour passer par un sommet ; une bosse brune et verte sur une droite en tirets ; en bas, quatre dessins de bandes : une droite verte bordée de deux courbes brunes qui s'écartent en bosses pour passer par des sommets (un en haut, un en bas ; puis un en haut, deux en bas), et un petit fuseau avec deux flèches opposées}

pas d'autres sommets de $\Sigma$ que $a$

pas d'autres sommets de $\Sigma$ que $b$ et $c$
\note{en bas à droite, deux légendes reliées par des accolades au dernier dessin : un fuseau brun autour de la droite verte, la rive supérieure passant par $a$, l'inférieure par $b$ et $c$, avec une flèche vers le haut et une vers le bas ; la première légende est au-dessus du dessin, la seconde au-dessous}

\page{156}
\[
\begin{array}{ll}
\rho_s = \sigma_2 \sigma_1 & \sigma_1 \\
\rho_f = \sigma_1 \sigma_0 & \sigma_1 = \rho_f^{\,\sigma_0} \\
\rho_s \rho_f = \sigma &
\end{array}
\]
\note{en haut à gauche ; la lecture des indices $s$ et $f$ de $\rho$ est incertaine}

\note{au-dessous, un croquis au crayon : deux droites $D$, $D'$ et une troisième $\Delta_s$ qui se coupent ; les orientations $\alpha$ (le long de la droite), $\beta$ (transverse) et $\omega$ (rotation), et les marques $\sigma_D$, $\sigma_{D'}$ aux extrémités d'une flèche transverse}

$\alpha, \beta, \omega$ associés

\note{en bas à gauche, un dessin en couleur : deux droites orange qui se croisent en un point $s$ ; une droite vert clair en tirets $D_0$, fléchée, passant par $s$, entre deux droites vert clair $D'$ (au-dessus, fléchée $\omega^{*}$) et $D$ (au-dessous) ; les orientations $\alpha^{*}$ (rotation) et $\beta^{*}$ (flèche brune transverse) ; le triangle marqué $T$ sous le croisement}

$\alpha^{*}, \beta^{*}, \omega^{*}$ \emph{non} associés

$\omega^{*} = - \alpha^{*} \wedge \beta^{*}$

$\omega = \alpha \wedge \beta$

\[
\begin{array}{lcl}
(D_0, \beta^{*}, \omega^{*}) & \overset{\sigma}{\longmapsto} & (D_0, -\beta^{*}, \omega^{*}) \\[2pt]
(D_0, \beta^{*}, \omega^{*}) & \overset{\sigma_0}{\longmapsto} & (D_0, \beta^{*}, -\omega^{*}) \\[2pt]
(D_0, \beta^{*}, \omega^{*}) & \overset{\sigma_2}{\longmapsto} & (D_0, -\beta^{*}, -\omega^{*})
\end{array}
\]
\note{au milieu de la page, un diagramme en étoile autour de $(D_0, \beta^{*}, \omega^{*})$, dont partent six flèches à talon : $\sigma$ vers le haut, $\sigma_0$, $\sigma_1$ et $\sigma_2$ vers la droite, et deux grandes flèches courbes $\rho_s$ et $\rho_f$ vers les dessins du haut ; nous le donnons ligne à ligne. Chaque but est suivi d'un petit dessin en couleur (le croisement orange, la droite en tirets déplacée d'un côté ou de l'autre, les flèches $\beta$, $\omega$ retournées). Le but de $\sigma_1$ est une parenthèse de deux dessins réunis par « ou », avec l'orientation $-\omega^{*}$ et le triangle $T$ ; les buts de $\rho_s$ et de $\rho_f$ sont de même des paires de dessins réunis par « ou », où apparaissent les triangles $T$ et $T_1$ (hachurés), et, pour $\rho_f$, une double flèche entre eux légendée « \ill{} » et la légende « premiers \ill{} de triangles »}

$\sigma_0$ \struck{change} change $\alpha$ en $-\alpha$, en gardant $\sigma_D, \sigma_{D'}, \beta$

$\sigma_2$ échange $\sigma_D$ et $\sigma_{D'}$, change $\beta$ en $-\beta$, garde $\alpha$

\page{157}
\note{en haut à gauche, une rangée de petits traits de couleur (essais de feutres)}

$X, \check{X}$ plans projectifs réels duaux, \struck{$T$ système \uncertain{« triangle »} de trois droites \ill{}} $T$ \add{\ill{} $X$}, $T'$ triangle dual dans $\check{X}$ \struck{et \uncertain{points} des}, \uncertain{dans} les orientations de $T$ et $T'$ se correspondant, de façon qu'il y ait $\omega'$ \uncertain{et} $\omega$ associés, d'où des orientations des \struck{\ill{}} droites d'\uncertain{origine} de $T$ et de $T'$, et de $X$ \uncertain{respect.}\ $\check{X}$ \uncertain{les} \ill{} \ill{} de $T$, \uncertain{resp.}\ $T'$, dans \uncertain{passer} \uncertain{un} \struck{pt} \add{\ill{}} de $X$ \ill{} $\check{X}$ et la droite correspondante de $\check{X}$ \ill{} $X$, les orientations y correspondant.
\note{texte en haut au milieu de la page, en colonne étroite, écrit sur l'impression d'un listing ; la première ligne biffée porte des guillemets au-dessus de « triangle » ; « $X$ » est ajouté au-dessus de la fin de la ligne biffée ; au-dessus de « pt » biffé, un mot illisible. Lecture d'ensemble très incertaine}

\note{à gauche, un grand dessin : trois droites numérotées $1$, $2$, $3$ (brune, noire, brun foncé) formant un triangle hachuré, orienté par une flèche de rotation ; en chacun des trois sommets, un point de couleur (bleu, vert, orange) entouré de deux arcs de la même couleur fléchés. À droite, le triangle dual : trois droites $1$, $2$, $3$ (bleu clair, orange, vert clair) aux sommets bruns, les arcs bruns et noirs fléchés, l'orientation $\omega'$. Au-dessous, deux croquis : une droite $1$ au point $2$, les vecteurs $\alpha$ (le long de la droite), $\beta$ (transverse) et la rotation $\omega$ ; une droite vert clair au point marqué, avec $\omega'$, $\beta'$, $\alpha'$}

$\alpha, \beta, \omega$ associés, mais $\alpha', \beta', \omega'$ \emph{non} \emph{associés}.

\note{en bas, deux dessins : deux droites orange qui se croisent, deux droites en tirets (bleue et verte) et une droite bleue pleine qui passent entre elles, les orientations $\alpha$, $\beta$, $\gamma$ marquées au croisement ; au-dessous, un petit croquis avec $\alpha$, $\beta$, $\gamma$ ; plus à droite, un segment brun épais entre deux points orange, avec en son milieu un point et les orientations $\beta'$, $\gamma'$, $\alpha'$}

\page{158}
\note{croquis au crayon, sans texte suivi : un rectangle partagé par une horizontale, les deux moitiés marquées $\alpha$ et $\beta$, un cercle au crayon en surimpression ; un rectangle traversé d'un disque aplati et de courbes qui se tressent ; un ovale ; au-dessous, deux rectangles superposés traversés de courbes tressées, les coins et les bords marqués $v'$, $u'$, $a$, $b'$, $u$, $v$, $u'$, $a'$ (?), $a$, $u$, $v$, avec deux ellipses aux niveaux $a$ et $b'$ ; à côté, un disque aux courbes entrelacées. En bas à gauche, un dessin en couleur : deux bandes vert clair qui se rejoignent en un point noir de croisement, une droite rouge qui les longe et une droite rouge verticale, des hachures au crayon et de petites barres noires}

\page{159}
\emph{1.1.84} \quad Complément de \uncertain{multiplication} sur les facettes exceptionnelles de $(\mathcal{X}, K)$
\note{la date est de sa main, en tête de la page, soulignée. Toute la page est barrée d'une grande croix au crayon vert ; elle est lue en entier}

a) \struck{\ill{}} Avec les notations de \uncertain{$i$}, on a
\[
\varphi^{-1}(X_1) = L \cup K_0 \qquad (\text{et } \ill{} = L\,!)
\]
\note{la lecture de l'indice de $X_1$ est incertaine}

b) La structure d'une facette exceptionnelle au voisinage de \struck{\ill{}} \add{\uncertain{deux}} arête\supplied{s} exceptionnelle\supplied{s} (en \uncertain{minime}, de la \add{double} \uncertain{paire} de facettes adjacentes exceptionnelles) est de \ill{} \ill{} \uncertain{sommet} la forme suivante :
\note{« deux » (lecture incertaine) est écrit au-dessus de la ligne, sur un mot biffé ; « double » est ajouté au-dessus de « paire »}

\note{dessin en couleur : deux facettes contiguës $F$ et $F'$, séparées et bordées par trois bandes verticales vert clair ; chacune est un cylindre dont les bords haut et bas sont des fuseaux rouges, traversé de courbes rouges verticales qui se croisent en huit ; des marques bleues aux coins de $F$ ; une droite horizontale (brune dans $F$, vert foncé dans $F'$) prolongée en pointillés bruns à gauche et à droite, marquée $D_i$ à droite, le point $a$ sur la bande commune ; des contours en tirets noirs, et les lettres $\delta$, $s$ dans $F$, $\delta'$, $s'$, $t'$ dans $F'$ ; les fuseaux de $F'$ se prolongent à droite en pointillés rouges}

\ill{} d'\uncertain{où} les « branches principales » \add{$\delta$} de $L$ dans une face exceptionnelle ($F$ ou $F'$ disons) \uncertain{par} un sommet \add{multiple} \struck{multiple}, dans les branches de $L$ en $s$ distinctes de \struck{\ill{}} $\delta$ \ill{} \uncertain{concourent} \ill{} un \uncertain{même} \uncertain{pt}.\ de $\delta F$ \uncertain{sur} $K_0$ (et par la loi de symétrie, \ill{} $= \ill{}$ \ill{} « un \uncertain{fuseau} » \ill{} $=$ multiplicité.
\note{« $\delta$ » est écrit au-dessus de « principales », rattaché par un trait ; « multiple » est réécrit au-dessus du mot biffé ; la parenthèse ouverte n'est pas refermée}

\uncertain{De} plus, je rappelle que la loi de symétrie s'étend à la multiplicité des pts \uncertain{doubles} sommets \struck{multiples} \add{de $L$} qui se trouvent sur les deux arêtes \uncertain{ordinaires} incidentes à $a$)

c) Pour \uncertain{toute} arête de $F$ distincte des \add{\uncertain{non exceptionnelles}} \add{de $F$} trois arêtes marquées, \struck{\ill{}} la \struck{\uncertain{bonne}} rive privilégiée est celle \uncertain{opposée} à cette face --- pour une arête \add{ordinaire} \uncertain{passant} \ill{} un sommet exceptionnel (arête \uncertain{\ill{}} \uncertain{marquée} ici), la \add{rive} \struck{arête} privilégiée est celle qui correspond à la face instable.
\note{l'addition « \uncertain{non exceptionnelles} de $F$ », écrite au-dessus de la ligne, est rattachée par une accolade ; « ordinaire » et « rive » sont écrits au-dessus de la ligne}

d) Pour avoir la position relative correspondant à une face exceptionnelle, il faut modifier la \uncertain{vraie} générisation \struck{\ill{}} (\uncertain{disons} \uncertain{générique} quand $\varphi(\partial F^{*})$ \ill{} : $F^{*}$ est la \uncertain{déduite} de $F$ par \uncertain{remplacement} \uncertain{en} les sommets critiques dans la rive \struck{exceptionnelle} privilégiée \uncertain{est} opposée à $F$) \uncertain{en} \struck{\uncertain{faisant} \ill{}} \add{\uncertain{définissant}} $F^{*}$ de façon \struck{\ill{}} \uncertain{similaire}, à partir de la « \ill{} \struck{\ill{}} modification obtenue en remplaçant \struck{des deux arêtes exceptionnelles des} \add{\ill{} bord d} exceptionnel » \struck{\ill{}} \add{de la} face envisagée (joignant \add{\uncertain{dans} \ill{}} les deux pts critiques principaux \struck{\ill{}} face) par un \struck{\ill{}} pointillé : l'intérieur de $F$ (cf.\ pointillé au stylo).
\note{colonne de droite ; nombreuses corrections interlinéaires, dont plusieurs illisibles ; le dernier mot biffé de la ligne « \uncertain{définissant} » est lu « faisant » sous toute réserve}

\page{160}
\note{page de dessins en couleur, sans texte : en haut à gauche, des droites rouges verticales numérotées en haut $1$ (surchargé), $1$, $2$, $3$, $4$, $5$ et en bas $1'$ (?), $1'$, $2'$, $3'$, $5'$, $4'$, coupées par deux bandes horizontales (vert foncé et vert clair) marquées $w$, $u$, qui se croisent au point $a$, une droite rouge fléchée au milieu ; des segments $x x'$ et $y y'$ entre les bandes, des hachures. À droite, trois disques : le premier bordé de vert clair, les points $a$, $a'$ au bord reliés par un fuseau rouge, le bord marqué $0', 1', 2', 3', 4', 5'$ et $0, 1, 2, 3, 4, 5$ avec des flèches ; le second, bordé de vert foncé, de même avec les points $b$, $b'$ ; le troisième réunit les deux, les bandes $a a'$ (vert clair) et $b b'$ (vert foncé) superposées, traversées de courbes rouges qui se croisent. En bas, un cylindre rouge (comme la facette de la page 159) ; un grand disque traversé de pseudo-droites rouges, le bord partagé en arcs vert clair et vert foncé ; un dessin semblable au troisième disque, prolongé à droite par un second disque, avec les points $a$, $a'$, $b$, $b'$, $x$, $x'$, $y$, $y'$ et des tirets ; en bas à gauche, deux bandes vert clair et vert foncé qui se croisent entre des droites rouges, les points $x$, $x'$, $y$, $y'$. En haut à droite, un disque au crayon traversé de courbes}

\end{document}
